\subsection{从谱积分到温度四次方}
在普朗克谱积分中令 $x=h\nu/(k_BT)$，则 $\dd\nu=(k_BT/h)\dd x$，所有温度因子可以提出积分号：
\[
 u=\frac{8\pi k_B^4T^4}{h^3c^3}
 \int_0^\infty\frac{x^3}{e^x-1}\,\dd x
 =\frac{8\pi^5k_B^4}{15h^3c^3}T^4.
\]
这里使用无量纲积分 $\int_0^\infty x^3/(e^x-1)\,\dd x=\pi^4/15$。结合 $j=cu/4$，得
$\sigma=2\pi^5k_B^4/(15h^3c^2)$。温度四次方既来自谱中的三次幂，也来自频率积分的一个尺度因子。

若球形天体半径为 $R$、距离为 $d\gg R$，忽略吸收且均匀辐射，则总光度 $L_*=4\pi R^2\sigma T^4$，探测器处辐照度
\[
 F=\frac{L_*}{4\pi d^2}=\sigma T^4\left(\frac Rd\right)^2.
\]
表面的辐出度与远处接收的辐照度同为 $\mathrm{W/m^2}$，却对应不同面积，不能因单位相同而直接相等。实际恒星大气、选择性吸收会改变谱形，本节只使用黑体近似。
